\documentclass[10pt,a4paper]{amsart}
\usepackage[margin=19mm]{geometry}
\usepackage{amsmath,amssymb,mathtools}
\usepackage[scr=rsfs,cal=euler]{mathalfa}
\usepackage{xfrac}
\usepackage[unicode,hidelinks,bookmarks=false]{hyperref}
\newcommand{\ie}{i.e.\ }
\newcommand{\cf}{cf.\ }
\newcommand{\resp}{respectively\ }
\newcommand{\Subsection}{\subsection}
\providecommand{\nobreakdash}{\nobreak\hbox{-}}
\setlength{\emergencystretch}{2em}
\multlinegap0pt
\usepackage{amssymb}
\usepackage{bm}
\usepackage{mathtools}
\usepackage{csquotes}
\usepackage[shortlabels]{enumitem}
\usepackage{xcolor}
\newtheorem{lemma}{Lemma}
\newtheorem{corollary}{Corollary}
\newcommand{\cC}{\mathcal{C}}
\newcommand{\cM}{\mathcal{M}}
\newcommand{\cR}{\mathcal{R}}
\DeclareMathOperator{\rc}{rc}
\DeclareMathOperator{\rr}{r}
\title{\bf Rank Contributions of Vertices in Rigidity Matroids of Clique Covered Graphs}
\author{Bill Jackson, Tibor Jord\'an, Soma Vill\'anyi}
\date{Source: arXiv:2607.26266v1 (CC BY 4.0)}
\begin{document}
\maketitle

\noindent\emph{Source and licence.} \bf Rank Contributions of Vertices in Rigidity Matroids of Clique Covered Graphs. Version arXiv:2607.26266v1 (CC BY 4.0), distributed by arXiv under the Creative Commons Attribution 4.0 International licence, http://creativecommons.org/licenses/by/4.0/, as recorded in the arXiv licence metadata for this exact version and shown on the abstract page https://arxiv.org/abs/2607.26266v1. Full licence notice: \texttt{licenses/arxiv-2607.26266-CC-BY-4.0.txt}.\par\smallskip

\noindent\textit{Editorial scope.} Counted unit: the article's lemma lem:canbecovered (source0.tex lines 1727--1732). The article's complete original proof of it is delivered with it (source0.tex lines 1734--1783, one \texttt{proof} environment of 50 lines). No mathematics is written, repaired or re-derived: the statement and the proof are the article's own lines, in the article's own notation, and no retained line is substituted or re-flowed.  Retained uncounted as the source-local context the proof needs: lines 661--663; lines 780--785; lines 1663--1667.  Nothing was omitted from any retained span, so the package carries no omission record.  No retained line contains a citation, so the wrapper carries no bibliography and no bibliography entry.  No retained line contains a figure environment, an image-inclusion command or drawing code, so no image is delivered and the package declares no asset dependency.  Editorial formatting only, in addition: an AMS \texttt{amsart} preamble that restates the article's own declarations and macros used by the retained lines, namely the environments \texttt{lemma}, \texttt{corollary} on separate counters, together with the standard packages those lines need.  The retained environments print in this document as Lemma 1, Corollary 1 in order of appearance, whereas the article prints the same items with the numbering of its own full document.  The complete native source of the article as distributed in the arXiv e-print for this exact version is bundled unchanged as \texttt{sources/206252/source0.tex}.
\begin{lemma} \label{lem:rc_sum}
For every subgraph $G$ of $K_n$, we have ${\displaystyle \rr(G)=\sum_{v\in V} \rc(G,v)}.$ 
\end{lemma} 

\begin{corollary}\label{coro:rc:geq:d}
Let $ G=(V, E)$ be a graph, $v\in V$ with $\deg_G(v)=k\geq d$, and
		$\cM\in\{\cR_d,\cC_{d-1}^{d-2}\}$.
        Suppose that $G+K(N_G(v))$ is $\mathcal{M}$-rigid but $G$ is not.
		Then $\rc(G,v)\geq d+\frac{1}{2}-\frac{1}{k}{\binom{d+1}{2}}.$
\end{corollary}

\begin{lemma}\label{lem:mainC5lemma}
    Let $G=(V,E)$ be a graph and $v\in V$ with $\deg_G(v)=k\geq 4$.
    Suppose that  $G-v+K(N_G(v))$ is rigid but $G$ is not. If each edge incident with $v$ belongs to a copy of $K_5$ in $G$, then
    $\rc(G,v)\geq 3$.
\end{lemma}

\begin{lemma}\label{lem:canbecovered}
    Let $H$ be a 5-connected $K_5$-covered graph  and let $w$ be a vertex of $H$. 
    Suppose that  $H[N_{H}(w)]$
    is clique bipartitionable.
    Then $H-w$ is rigid.
\end{lemma}

\begin{proof}
    Suppose, for a contradiction, that the lemma is false. Let $H$ be a counterexample with as few vertices as possible and, subject to this condition, as many edges as possible.
    Let $G=(V,E)=H-w$. Then $G$ is not rigid and
    the maximality of $|E(H)|$ implies that
    \begin{equation}\label{eq:rigcluster}
        \mbox{each rigid cluster of $G$ of size at least five induces a clique in $G$.}
    \end{equation}
    For each vertex $v\in V$ let $H_v=H-v+K(N_{H}(v))$. Then $H_v$ is a 5-connected $K_5$-covered graph, for which $H_v[N_{H_v}(w)]$ is clique bipartitionable.
    Thus the minimality  of $|V(H)|$ implies that  $H_v-w=G-v+K(N_G(v))$ is rigid.  
    We will complete the proof by calculating the rank contributions of the vertices of $G$.
    
    We first consider a vertex $v\in V-N_{H}(w)$.
    Since $H$ is a $K_5$-covered graph, 
    each edge incident with $v$ belongs to a copy of $K_5$ in $G$.
    Lemma \ref{lem:mainC5lemma} now implies that 
    \begin{equation}
        \label{out}
        \mbox{$\rc(G,v)\geq 3$   for each $v\in V-N_{H}(w)$.}
    \end{equation}
    
It remains to consider a vertex  $v\in N_H(w)$. 
    Since $H[N_{H}(w)]$
    is clique bipartitionable, and $H$ is a $5$-connected $K_5$-covered graph,
    there is a bipartition 
    $N_H(w)=V_1\cup V_2$, such that $|V_1|,|V_2|\geq 2$,
    and $G[V_1]$ and $G[V_2]$ are both cliques.
    By relabelling, if necessary, we may assume that $v\in V_1$.
    Since $G-v+K(N_G(v))$ is rigid, $G$ is not rigid, and $\deg_G(v)\geq 4$,
    it follows that
    $K(N_G(v))\not\subseteq G$. 
    Let $V_1'$ be a maximal subset of $V$ with $V_1\subseteq V_1'$ for which
    $G[V_1']$ is a clique and let $x\in N_G(v)-V_1'$.
    Since $H$ is a $K_5$-covered graph, 
    there exists a set $V'\subseteq V$ such that $V'$ induces a maximal clique in $G$,  $|V'|\geq 4$, and  $v,x\in V'$.
    We can now use \eqref{eq:rigcluster} to deduce that $|V_1\cap V'|\leq 2$.
    Thus $|V'-V_1|\geq 2$, and hence 
    $\deg_G(v)\geq |V_1|+1$. 
By applying
    Corollary \ref{coro:rc:geq:d}
    and using that $|V_1|\geq 2$ we can now deduce that
$\rc(G,v)\geq \left(3+\frac{1}{2}-\frac{6}{|V_1|+1}\right)$ and hence
$$\sum_{v\in V_1}\rc(G,v)\geq |V_1|\left(3+\frac{1}{2}-\frac{6}{|V_1|+1}\right)\geq 3|V_1|+\left(\frac{|V_1|}{2}-\frac{6|V_1|}{|V_1|+1}\right) \geq 3|V_1|-3.$$
    By symmetry, we also have $\sum_{v\in V_2}\rc(G,v)\geq 3|V_2|-3.$
    We can now apply Lemma \ref{lem:rc_sum} and (\ref{out}) to obtain
    \begin{align*}
        \rr(G)=\sum _{v\in V}\rc(G,v)&= \sum_{v\in V\setminus (V_1\cup V_2)}\rc(G,v)+\sum_{v\in V_1}\rc(G,v)+\sum_{v\in V_2}\rc(G,v)\\
        &\geq  3|V\setminus (V_1\cup V_2)| + \left(3|V_1|-3\right)+\left(3|V_2|-3\right)= 3|V|-6.
    \end{align*}
    This contradicts the assumption that $G$ is not rigid.
\end{proof}

\end{document}
